\section{Exact Risk and Scaling Law}

For $K\sim\mathrm{Binomial}(B,q)$ define
\begin{equation}
h_B(q):=\E\left[\frac{\1\{K>0\}}{K}\right].
\label{eq:h-def}
\end{equation}

\begin{theorem}[Exact finite-block risk]
\label{thm:exact-risk}
For every $B,M\ge1$ and every target with
$\sum_j s_j<\infty$, the estimator in Equation~\eqref{eq:estimator} satisfies
\begin{align}
\E\risk(\widehat\theta,\theta)
={}&\underbrace{\sum_{j>M}s_j}_{\text{model approximation}}
+\underbrace{\sum_{j\le M}s_j(1-q_j)^B}_{\text{unseen-mode bias}}
\nonumber\\
&+\underbrace{\sigma^2\sum_{j\le M}
\frac{\lambda_j}{\tau_j}h_B(q_j)}_{\text{label-noise variance}}.
\label{eq:exact-risk}
\end{align}
\end{theorem}

\begin{proof}
If $K_j=0$, coordinate $j\le M$ is estimated by zero and contributes
$s_j$. If $K_j=k>0$, the signed average is unbiased and has conditional
variance $\sigma^2/(\tau_jk)$. Since $K_j\sim\mathrm{Binomial}(B,q_j)$,
averaging these cases and adding the unavoidable $j>M$ tail gives
Equation~\eqref{eq:exact-risk}.
\end{proof}

The variance term is governed by a simple reciprocal-count estimate.

\begin{lemma}[Reciprocal binomial count]
\label{lem:binomial}
There are universal constants $0<c<C<\infty$ such that, for all $B\ge1$ and
$q\in(0,1]$,
\begin{equation}
c\min\{Bq^2,B^{-1}\}
\le qh_B(q)
\le C\min\{Bq^2,B^{-1}\}.
\label{eq:binomial-bound}
\end{equation}
\end{lemma}

The proof is given in the technical supplement. We now state the principal scaling
law. The notation $u_j\asympconst v_j$ means that their ratio is bounded above
and below by positive constants independent of $j$, $B$, and $M$.

\begin{theorem}[Persistent spectral scaling law]
\label{thm:scaling}
Assume
\begin{equation}
\lambda_j\asympconst j^{-a},
\qquad
\tau_j\asympconst j^r,
\qquad
s_j\asympconst j^{-b},
\label{eq:power-laws}
\end{equation}
where $a>1$, $r\ge0$, and $b>1$. Put $p=a+r$. Then, uniformly for
$B,M\ge2$,
\begin{equation}
\boxed{
\E\risk(\widehat\theta,\theta)
\asympconst
M^{-(b-1)}
+B^{-(b-1)/p}
+\sigma^2\frac{\min\{M,B^{1/p}\}}{B}.}
\label{eq:main-scaling}
\end{equation}
The hidden constants depend only on the comparison constants in
Equation~\eqref{eq:power-laws} and on $a,b,r$.
\end{theorem}

\paragraph{Interpretation.}
The innovation probabilities satisfy $q_j\asymp j^{-p}$. A mode has a
constant probability of appearing once $Bq_j\gtrsim1$, giving the spectral
resolution
\begin{equation}
J_B\asymp B^{1/p}=B^{1/(a+r)}.
\label{eq:JB}
\end{equation}
The first two terms in Equation~\eqref{eq:main-scaling} are equivalently
$\Theta(\min\{M,J_B\}^{-(b-1)})$. The third term counts the number of
well-observed modes, up to the same resolution, divided by $B$.

\begin{corollary}[When persistence changes the exponent]
\label{cor:exponent}
In the noiseless case, or for fixed $\sigma^2=O(1)$ in the hard-target regime
$1<b<a+r$ where the bias term dominates the variance term,
\begin{equation}
\E\risk(\widehat\theta,\theta)
\asympconst M^{-(b-1)}+B^{-(b-1)/(a+r)}.
\label{eq:clean-law}
\end{equation}
Uniform persistence has $r=0$ and preserves the i.i.d. exponent
$(b-1)/a$. If weaker modes persist longer with $r>0$, the data exponent becomes
$(b-1)/(a+r)$.
\end{corollary}

\begin{corollary}[Noise regimes and model choice]
\label{cor:model-choice}
Let $\sigma>0$ be fixed. If $b<a+r$, taking
$M\asymp B^{1/(a+r)}$ achieves risk
$\Theta(B^{-(b-1)/(a+r)})$. If $b>a+r$, the rate-optimal model size is
\begin{equation}
M_\star\asymp(B/\sigma^2)^{1/b},
\label{eq:optimal-M}
\end{equation}
and the resulting dominant risk is
$\Theta((\sigma^2/B)^{(b-1)/b})$.
\end{corollary}

